% preface-body.tex -- the preface, with no preamble, for \input into a volume.
% GENERATED from preface.tex (which remains the standalone build).

\frontmatter

\chapter*{Preface}
\addcontentsline{toc}{chapter}{Preface}

This book series forms part of the research program on \emph{generative relational being}. Before undertaking the series, we set out a diagnosis of the contemporary condition in the foundational paper \emph{Responding to the Crises of Symbol and Subject in Modernity: Towards a Generative Nomos of Relational Presence}. The argument of that paper begins with the constitution of the subject: the subject is constituted within relations, through symbolic mediation by language, norms, law, and institutions. The symbolic order provides the field in which the subject emerges, takes form, attains presence, and undergoes transformation, while the ontological ground of the subject remains relational and exceeds every symbolic representation. In modernity, the rapid proliferation of symbolic production, together with the steady weakening of the symbolic order's mediating capacity, has rendered the intrinsic gap between the Real and the Symbolic visible to an unprecedented degree.

This gap manifests today in a series of concrete phenomena. Grand narratives, among them justice, progress, and divine order, have lost their persuasive force; public trust in institutions continues to drain away, shared norms fragment, and moral and cultural life slides into relativism. Digital media and artificial intelligence have made the production of signs nearly costless, so that symbols are devalued through abundance and empty repetition; beneath the appearance of novelty lies structural stagnation, and meaning enters a state of inflation. The intersubjective space is increasingly occupied by interfaces, algorithms, and bureaucratic systems; the I--Thou relation degrades step by step into an I--It operation, and genuine encounter and mutual recognition grow ever thinner. What the symbolic order has long repressed, the Real, returns in traumatic form: conspiracy narratives attempt to give shape to diffuse anxiety, fundamentalist violence erupts intermittently, and a free-floating dread pervades conditions of material abundance. Where relational trust collapses, individuals and communities turn to law for order, for meaning, and even for redemption, so that the juridical system bears expectations far beyond what its structure can sustain and justice falls into overload. Algorithmic recommendation systems form an epistemic loop in which the subject's knowledge of the world reduces to the repeated confirmation of its own mediated and pre-shaped image; affect, in turn, is simulated and recirculated in interactions with constructed pseudo-subjects and gradually loses its grounding in the Real.

Placed in historical depth, these phenomena disclose three epochs in the position of the big Other. In the first epoch, gods and natural forces occupied the position of the big Other, grounding order and signification; humanity constructed its existence through cosmological and sacred orders. The Enlightenment opened the second epoch: symbolic reason replaced external authority, and rationality, legality, and institutional norms formed a new regime of mediation, with law occupying a privileged position as the practical embodiment of the big Other. The present constitutes the third epoch: the credibility and mediating capacity of symbolic structures are continually destabilized, and as symbolic authority deteriorates the subject seeks to approach, and even to occupy, the position of the big Other, so that cycles of integration, struggle, and conflict recur. Along the extension of this genealogy, a more fluid state of mediation becomes conceivable, resonant with the Daoist ideal of \emph{wuwei} and of harmony between heaven and humanity: the subject moves from being a product of structure to being an active participant in its generativity, and existence turns once more toward relationality itself.

From this follows the central diagnosis of the foundational paper: the crisis of subjectivity and the symbolic crisis are mutually generative, each serving as cause and consequence of the other, and the widening gap between them constitutes the core predicament of modern existence. The crisis originates in the gap between the Real and the Symbolic; any response must therefore address the relation between being and its mediation, since efforts confined to repairing symbolic structures remain at the level of symptoms. Our response returns to the ontological ground of existence. The first element is relational ontology: the basic units of being are relations, relation precedes substance, and the self is an emergent formation within a relational field. The second element is generative presence: the epistemic and phenomenological process through which relational being is enacted in lived experience and through which meaning comes forth. In the domain of law and justice, this entails a double movement, juridical and extra-juridical: the former reorients law toward relational being and toward the ongoing processes through which its authority and meaning are produced and sustained; the latter cultivates the formation and expressive realization of generative presence within relational existence. The diagnosis thus opens onto a direction: towards relational being.

Moving towards relational being requires enabling conditions. Generative relational being therefore constitutes an infrastructural research program, one that investigates the conditions under which relational being can be understood, enacted, generated, sustained, and carried across contexts. Infrastructure here designates a set of enabling conditions, conceptual, theoretical, procedural, social, and ethical. The program organizes these conditions into seven interrelated infrastructural domains: a theoretical infrastructure, providing formal and modelable frameworks for articulating relational structures and generative processes; a conceptual and ontological understanding infrastructure, clarifying the core principles of relation, being, presence, structure, and generation; a co-presence infrastructure, supporting the experiential conditions through which relations are perceived and enacted across distance and mediation; a decision infrastructure, enabling relational action and governance to unfold through accountable, process-oriented, and evidence-informed practices; a co-creative infrastructure, sustaining the generative emergence of new relations, meanings, and structures through shared practices; an educational infrastructure, cultivating and transmitting relational capacities across individuals, communities, and generations; and an ethical and care infrastructure, carrying and reconstituting relations under conditions of vulnerability, asymmetry, and dependence.

This book series is a research instance within that program: intimate relationships.

We take intimacy as our point of entry because it constitutes a relatively small yet structurally complete relational system. The problems that the seven infrastructures address, among them recognition and misrecognition, co-presence and symbolic mediation, joint decision and asymmetries of power, the generation and alienation of value, care and vulnerability, eros and justice, all appear in intimate relationships in highly condensed form. Intimacy is comparable to a tabletop experiment in physics: small in scale, yet sufficient to make manifest the laws that govern much larger social systems. Further, if the substance of the symbolic crisis lies in the separation between lived relation and its symbolic representation, then intimate relationships are the very place where this separation is most immediately experienced and where it stands the best chance of being generatively repaired. The series examines the relational processes unfolding between two subjects in order to discover the generative mechanisms and the claims of justice that run through all relational systems; the findings obtained here will in turn flow back into the program as a whole.

This work has grown at the intersection of several theoretical traditions. Since the provenance of these ideas itself belongs to the problematic of the series, the order in which the influences entered deserves an honest account. At the outset of the series, the resources already in place included: the feminist revaluation of care, dependence, and vulnerability; jurisprudence and the theory of justice; the Kantian normative framework of self-legislation and autonomy, together with the Hegelian dialectic of recognition; Saussurean structural linguistics, which shows that meaning is generated within systems of difference; Lacanian psychoanalysis, on the subject constituted in the desire of the Other and misrecognizing itself in the mirror; the empirical findings of developmental psychology and neuroscience on joint attention, attachment, and the formation of the subject; Miranda Fricker's theory of epistemic injustice; and the formal and systematic methods of theoretical physics (the purely relational ontology suggested by field theory and symmetry), engineering, and sociology. These formed the initial theoretical ground of the series.

A second set of resources, of equal importance, entered and reshaped the problematic of the series only in the course of writing, above all through scholarly exchange with Ron Eglash. The framework of generative justice shows that the crux of justice lies in whether value can flow back to its creators at the very site of its generation and circulate within relational networks free from alienated extraction; the redistribution of already-constituted value is a derivative form of this more fundamental question. Along this thread, political economy, Marx and post-Marxist philosophy, and anthropology successively entered the theoretical horizon of the series. They extended the cluster of questions concerning how value is created, circulated, and extracted from the domains of ecology and labor to intimate relationships, the most immediate site of value generation: does the value created by love, attention, and care circulate within the relation, or is it extracted in one direction?

This turn also has a personal origin. The author's partner was trained in political economy and cultural studies. It was she who made the author aware of a gap in his own theoretical apparatus; and it was out of the motivation to understand her truly, and to meet her within a shared vocabulary, that the author undertook a systematic study of the theories and frameworks of political economy and cultural studies. The problematic of the latter half of the series is saturated with the traces of this encounter. We record this explicitly here because it confirms one of the central theses of the series: theoretical work is itself a relational practice, and thought is here the effort of one subject to draw near to another.

These traditions are by no means always in harmony; at times they stand in mutual tension. The series works precisely within that tension, seeking to provide for relational being a theoretical apparatus that combines philosophical argument with a formal skeleton.

A word is also owed to the reader concerning the form of the series. The series takes the form of a collection of independent papers gathered into books. This form has its scholarly tradition: Lacan's \emph{\'Ecrits} is a compilation of papers written over many years, held together by a single sustained problematic; Benjamin's essay collections take the form of a constellation, in which the pieces illuminate one another and meaning is generated in the intervals between them. The series follows this practice. The volumes enter from different sides, including the normativity of self-legislation in intimate relationships, joint attention and the formation of the subject, epistemic injustice in the situation of the marriage proposal, and the circulation and alienation of value within relationships; each paper retains the integrity of an independent argument, while cross-references and a shared problematic weave the papers together. This form is itself isomorphic with the central thesis of the series: meaning is generated within the relations of a system of differences. A single thesis runs through all the volumes: intimate relationships are the very place where the question of justice is raised first and raised most sharply. How we treat those closest to us is our most honest answer to the question of justice.

Concerning the methodological position and the expectations of the series, a further word is owed to the reader. The main body of the series belongs to theoretical research and philosophical reflection; its methods are primarily those of traditionalist social science methodology and philosophical argument: the clarification of concepts, the construction of arguments, the critical synthesis of traditions, and, on this basis, attempts at formalization. This position means that many of the propositions of the series are, in character, testable points of departure. From this we hold a twofold expectation. First, we expect positivist methodology to carry this work forward: to translate the concepts and propositions proposed here into operational research questions, to test and revise them in the domains of experience and practice, and to discover new problems that theoretical reflection alone cannot reach; the study of intimate relationships possesses deep empirical traditions in developmental psychology, sociology, and neuroscience, and the series wishes to be an interlocutor of, and a questioner for, those traditions. Second, we expect the questions raised by the series to offer a valuable point of reference at the level of social governance: if intimate relationships are indeed the place where questions of justice first take visible form, then the insights obtained here concerning recognition, trust, the circulation of value, and care should assist reflection on institutional design, public policy, and communal life at larger scales. This twofold expectation is at the same time a self-limitation: what the series hands over are questions and frameworks, and their empirical validity and practical utility await the answers of those who come after, each in their own field.

Finally, concerning the openness of the series. The source of the series will be made public; the public version removes the acknowledgments, the dedication, and other private passages, so that it stands as a purely philosophical work, a set of propositions and arguments that others may read, cite, modify, contest, continue, and even overturn. This decision follows directly from the central thesis of the series. If relational being is indeed generative, if value emerges and circulates within networks of relation, then a work arguing that the subject comes to completion within relations ought to exist in a manner commensurate with its own thesis. The themes treated in the series, among them how the subject is constituted in the desire of the Other, how value emerges in care and recognition, and how justice is claimed under asymmetries of power, all exceed the exhaustive grasp of any single mind. This is a consequence of the theory of the series itself: the perspective of any author is necessarily partial, situated, and incomplete. To acknowledge this, and to open the gaps to the supplementation, questioning, and rewriting of others, is the mode of existence commensurate with the subject matter of the series. This stance answers directly to generative justice: open sourcing brings the series into the circulation of value, and the new value generated by later modifications and derivations should return to this ever-expanding community of creators; a work that closes itself off, with its author monopolizing all rights of interpretation, falls into the extractive structure of the domain of knowledge and congeals into private property a meaning that ought to circulate. What we hand over is the present cross-section of a relational system still in generation. No one should pretend to possess the unlimited reason that would master the subject of this book; we extend an invitation to relation.

Throughout the writing of this series there has been one concrete person, one concrete relationship, serving as the silent reference for the refinement of every concept. May this private reference at last overflow the bounds of two people, and benefit many more who dwell within intimate relationships; may they generate one another within relation, and be just to one another.

\begin{flushright}
Tokyo\\
June 2026
\end{flushright}
